\par\noindent\textbf{(1)} 极板自由电荷满足 \(\dot Q=I(t)\)，板间电位移 \(D_z=Q/(\pi a^2)\)。因此 \(\partial_tD_z=I(t)/(\pi a^2)\)，穿过半径 \(r\le a\) 圆盘的位移电流为 \(I_d(r)=I(t)r^2/a^2\)，在 \(r=a\) 处等于导线传导电流。
\par\noindent\textbf{(2)} 对半径 \(r<a\) 的同轴圆环，积分式 \(\oint\boldsymbol H\cdot d\boldsymbol l=d/dt\int\boldsymbol D\cdot d\boldsymbol S\) 给 \(2\pi rH_\phi=I(t)r^2/a^2\)。故 \(H_\phi(r)=I(t)r/(2\pi a^2)\)，磁场在轴心为零并随 \(r\) 线性增大。
\par\noindent\textbf{(3)} 板间电压 \(V=Qd/(\epsilon\pi a^2)\)，储能 \(U=Q^2d/(2\epsilon\pi a^2)=CV^2/2\)。时间导数 \(dU/dt=(Q/C)\dot Q=VI\)，与理想无损电容的输入电功一致；电磁场储能由边缘Poynting通量进入间隙。
\par\noindent\textbf{(4)} 同一闭合导线环路可张成穿过导线的曲面，也可张成穿过电容间隙的曲面。仅用 \(\oint Hdl=I_{\mathrm{cond}}\) 时两个曲面分别包住 \(I\) 与零，造成矛盾；引入 \(\partial_t\boldsymbol D\) 才使结果与曲面选择无关。
